\documentclass[a5paper,pagesize,10pt,bibtotoc,pointlessnumbers,
normalheadings,DIV=9,twoside=false]{scrbook}

\usepackage{graphicx}
\usepackage[babel,german=guillemets]{csquotes}

\usepackage[parfill]{parskip}

\linespread{1.05} 

\newcommand{\q}[1]{>>\textit{#1}<<}

\title{A book title}   
\author{Author Name} 
\date{\today} 

\begin{document}


%=========================================
\begin{titlepage}
        \centering{
            \includegraphics[width=0.9\linewidth]{wide_dark_nobg.png}
        }\\
        
		\centering{
			{\Huge{
            New member handbook}
            }
		}\\
			
		\vspace{10mm}
		
		\vspace{\fill}
		\centering \large{Nov. 2026}
\end{titlepage}

%=========================================

\newpage{}
\thispagestyle {empty}

\section{Important contacts}

For \textit{urgent} matters (e.g. a broken door or window, leaks), contact South Block:

\begin{itemize}
    \item Email: teia.hughes@waspsstudios.org.uk
    \item Office phone: 0141 271 4700
    \item Direct phone: 0141 553 5896 (ext. 1025)
    \item \textbf{Emergency out of hours: 0141 553 5890}
\end{itemize}

For \textit{non-urgent} matters:

\begin{itemize}
    \item Email the board: glasgowmakercollective@proton.me
    \item Ping the board on discord/irc with @board
\end{itemize}

\section{Opening the hackspace}\label{opening}

If you are in the space, you can optionally `open' it by flipping the red switch next to the door.
This will post a message to the \#space channel on Discord, announcing that the space is open to visitors.

If the space is open, you are responsible for managing guests, which may ask for access to the space. See section \ref{guests} for more information.

\section{Closing the hackspace}

If you are the last person in the space, you are responsible for `closing' the space when you leave. 

A non-complete checklist (also by the door) of things to do:
\begin{itemize}
    \item turn off:
    \begin{itemize}
        \item 3d printer
        \item soldering station
        \item beeg plotter
        \item lights
    \end{itemize}
    \item close the windows
    \item flip the space status switch to `off'
    \item ensure the space key is in the keybox outside the door
\end{itemize}

\section{Incidents}

If you have an accident, it's important to tell the board (for insurance purposes).
You can email the board 
\\(glasgowmakercollective@proton.me), 
or tell a board member in person.

\section{Non-member guests} \label{guests}

You are welcome to invite guests to the space!
They can use the space, including equipment.
It's important to note that \textbf{you} are then responsible for that person, including their safety, and damage they might do to the space or shared equipment.

\subsection{Other guests}

When the space is open (see section \ref{opening}), the public is invited to visit the space.
If someone does decide to visit, \textbf{all} members (and especially board) present are responsible for those people.
This includes watching out for their safety, and reporting incidents (e.g. injuries) to the board.

\section{Receiving mail}

You can have mail sent to the space. 

Large parcels are delivered to a nook behind sliding doors behind the door immediately to the right as you enter the main door to South Block.
Smaller letters are delivered to pigeon holes on the first floor of the building, opposite the shared kitchen.

The address to use is:

\begin{displayquote}
put your name here \\
Glasgow Hackerspace, Studio 221\\
60-64 Osborne Street\\
Glasgow\\
G1 5QH\\
\end{displayquote}

\section{Caring for the hackspace}

The hackspace is maintained by all of its users - that includes all members, as well as your guests!

To keep the hackerspace in a comfortable state for everyone, try to follow the `scout' rule: \textbf{always leave the space as you found it, or better}. For example, if you use a mug or cup, wash it afterwards!

Things you can do to help maintain the space include:
\begin{itemize}
    \item tidy workspaces and tables
    \item wipe down tables with cleaning spray
    \item returning chairs to the tables
    \item stacking stools
    \item hoover the floor 
    \item re-organise the electronics workbench
    \item restocking mate in the fridge
\end{itemize}

\section{Improving the hackspace}

The hackspace is an ongoing project by all members.
As such, you can improve things! 

We run a monthly meeting for the space, currently on the first Saturday of the month. 
Agendas can be found on the wiki, and are usually posted as a message on Discord at some point.
Everyone is welcome to add to the agenda of these meetings!

For smaller things (think adding lights, decorations, or anything else non-permanent), don't be afraid to just do it.
This ethos extends to running events such as meetups around a topic.
You can always ask the membership or board if you are unsure.

\subsection{Decision making}

Bigger decisions that affect everyone should be made with the membership, usually by vote.
For things that relate to the company itself, this will be done at either an AGM or an EGM.
For things relating to the space, this will usually be done in an ad-hoc way.

\newpage
\end{document}

